% !PRINTED 188
\weberpage{188}
%<<CONT
\[
\begin{array}{r@{}c@{}r@{}r@{}r}
\displaystyle\frac{\partial^2\sigma}{\partial u^2} &{}={}& -\dfrac{g_2}{12}u^3 &{}- 12\,g_3\dfrac{u^5}{2^4\cdot3\cdot5} &{}- \dfrac{3}{8}g_2^2\dfrac{u^7}{2^3\cdot3\cdot5\cdot7},\\[2mm]
\displaystyle\frac{\partial\sigma}{\partial g_2} &{}={}& &{}- \dfrac{u^5}{2^4\cdot3\cdot5} &\\[2mm]
\displaystyle\frac{\partial\sigma}{\partial g_3} &{}={}& & &{}- \dfrac{u^7}{2^3\cdot3\cdot5\cdot7},\\[2mm]
\displaystyle\frac{g_2}{12}u^2\sigma &{}={}& \dfrac{g_2}{12}u^3 & &{}-\dfrac{7}{24}g_2^2\dfrac{u^7}{2^3\cdot3\cdot5\cdot7},
\end{array}
\]

und wenn man dies in (12) einsetzt und die Koeffizienten von $u^5$
und $u^7$ gleich Null setzt, so ergibt sich:

\[
h_2 = -12\,g_3, \qquad h_3 = -\tfrac{2}{3}g_2^2,
\]

und wir erhalten die Differentialgleichung:

\weq{13}{\dfrac{\partial^2\sigma}{\partial u^2} - 12\,g_3\dfrac{\partial\sigma}{\partial g_2} - \dfrac{2}{3}g_2^2\dfrac{\partial\sigma}{\partial g_3} + \dfrac{g_2}{12}u^2\sigma = 0.}

Setzt man die Reihenentwickelung für $\sigma$ mit unbestimmten
Koeffizienten an:

\weq{14}{\sigma = \Sigma A_\nu\dfrac{u^{2\nu+1}}{(2\nu+1)!},}

so erhält man aus (13) die Rekursionsformel:

% [sic] The source prints g_2^2 in the final term of (15); coefficient extraction from (13), and the homogeneous recurrence (17), require g_2.
\weq{15}{A_\nu = 12\,g_3\dfrac{\partial A_{\nu-1}}{\partial g_2} + \dfrac{2}{3}g_2^2\dfrac{\partial A_{\nu-1}}{\partial g_3} - \dfrac{(2\nu-1)(2\nu-2)}{12}g_2^2A_{\nu-2},}

woraus zu schließen ist, daß die $A_\nu$ \gesperrt{ganze rationale Funktionen}
von $g_2$ und $g_3$ sind, deren Koeffizienten rationale Zahlen sind,
die nur Potenzen von 2 und von 3 im Nenner haben können. Die Reihe
(14) besteht daher aus Gliedern von der Form

\[
a_{m,n}\left(\tfrac12 g_2\right)^m(2g_3)^n\dfrac{u^{2\nu+1}}{(2\nu+1)!},
\]

und aus der Homogenität der Funktion $\sigma$ [§ 37, (8), § 46,
(13)] folgt:

\weq{16}{2m+3n=\nu,}

für die Koeffizienten $a_{m,n}$ ergibt sich dann aus (15):

\weq{17}{\begin{aligned}
a_{m,n} &= 3(m+1)a_{m+1,\,n-1} + \dfrac{16}{3}(n+1)a_{m-2,\,n+1}\\
&\quad{}- \dfrac{1}{3}(2m+3n-1)(4m+6n-1)a_{m-1,\,n},
\end{aligned}}

worin $a_{00}=1$ und die $a$, bei denen einer der beiden Indizes
negativ ist, gleich Null zu setzen sind.
